% ══════════════════════════════════════════════════════════════
%  Chapter 11: Storage Layer --- Persistence for Blocks and State
% ══════════════════════════════════════════════════════════════
\chapter{Storage Layer}
\label{ch:storage}

\begin{learnbox}
\begin{itemize}
  \item Use the Sled embedded database for persistent block and chain state storage
  \item Understand how the \code{BlockchainFileSystem} abstracts low-level storage operations
  \item Implement efficient persistence and retrieval of blocks, chain state, and UTXO data
  \item Explain the trade-offs in the storage design and why Sled was chosen
\end{itemize}
\end{learnbox}

\lettrine{W}{ithout} \index{durable storage}durable storage, a \index{node}node would lose all \index{blockchain}blockchain \index{history}history every time it restarts --- it would have to re-download and re-validate the entire \index{chain}chain from \index{peer}peers. The \index{storage layer}storage layer also enforces a critical \index{invariant}invariant: the \index{chain tip}tip \index{hash}hash and the \index{block}blocks tree must always be \index{consistency}consistent. If the \index{node}node crashes halfway through writing a new \index{block}block, the \index{database}database must either contain both the \index{block}block and the updated \index{chain tip}tip, or neither. This is why \index{atomicity}atomicity is the central design concern of this chapter.

% ──────────────────────────────────────────────────────────────
\section{What ``Storage'' Means in Our Rust Bitcoin Implementation}
\label{sec:ch11-what-storage-means}

In this codebase, ``storage'' is not a separate service; it is an embedded key-value database (\textbf{\index{sled}sled}) with a small set of conventions.

% ──────────────────────────────────────────────────────────────
\subsection{What Is Sled?}
\label{sec:ch11-what-is-sled}

\index{sled}Sled is a pure-Rust embedded database --- think of it as a persistent \code{BTreeMap<Vec<u8>, Vec<u8>>} backed by a log-structured merge tree on disk. It requires no external process, no TCP connection, and no configuration file: you open a directory path and get a \code{sled::Db} handle. We chose it because it gives us atomic batch writes (critical for ``insert \index{block}block + move \index{chain tip}tip'' in one operation) and compiles on every platform Rust targets. If you have used RocksDB or LevelDB before, \index{sled}sled fills the same niche with a Rust-native API.

Our conventions on top of \index{sled}sled:

\begin{itemize}
  \item a \textbf{\index{block}blocks tree} that maps \code{\index{block}block\_hash -> serialized \index{block}Block bytes}
  \item a stable \textbf{\index{chain tip}tip key} (\code{"tip\_block\_hash"}) that points at the canonical \index{chain tip}tip hash
  \item atomic updates for ``insert \index{block}block + move \index{chain tip}tip'' (\index{sled}sled transactions)
\end{itemize}

Think of sled as a persistent in-memory data structure (like a \code{BTreeMap}) that writes changes to disk atomically. If the process crashes, the database either contains both the block and the updated tip, or neither --- it never gets stuck in an inconsistent state.

\begin{notebox}[Debugging Tip]
When debugging storage issues, Sled's \code{Db::export()} and \code{Db::import()} methods can dump the entire database contents for inspection. This is invaluable during development.
\end{notebox}

The key methods are \code{BlockchainFileSystem::create\_blockchain} and \code{open\_blockchain} (lifecycle), \code{get\_tip\_hash} and \code{get\_last\_block} (reads), and \code{update\_blocks\_tree} (the internal atomic write that inserts a block and advances the tip in a single sled transaction).

% ──────────────────────────────────────────────────────────────
\section{Sled Layout: The Minimum Schema}
\label{sec:ch11-sled-layout}

\begin{minted}[fontsize=\footnotesize,linenos=false]{text}
sled::Db at ./<TREE_DIR>/
  |
  +- Tree "<BLOCKS_TREE>"  (default: "blocks1")
       |
       |-- key: "<block_hash_string>"
       |    -> value: serialized Block (bytes)
       |
       +-- key: "tip_block_hash"
            -> value: "<block_hash_string>" (bytes)
\end{minted}

This chapter is about making those two invariants true:

\begin{itemize}
  \item if a \index{block}block exists, it can be fetched by hash
  \item \index{chain tip}tip always points at a \index{block}block that exists (or a known ``empty'' sentinel)
\end{itemize}

% ──────────────────────────────────────────────────────────────
\section{Initialization: Creating or Opening the DB}
\label{sec:ch11-initialization}

The storage location is chosen by environment variables:

\begin{itemize}
  \item \code{TREE\_DIR} (directory on disk; default \code{"data1"})
  \item \code{BLOCKS\_TREE} (sled tree name; default \code{"blocks1"})
\end{itemize}

\subsection{Diagram: Config $\to$ Path $\to$ DB + Tree}
\label{sec:ch11-config-diagram}

\begin{minted}[fontsize=\footnotesize,linenos=false]{text}
TREE_DIR="data1"      BLOCKS_TREE="blocks1"
      |                    |
      v                    v
./data1/  (sled database)  Tree("blocks1")
\end{minted}

\subsection{Type Definition and Blockchain Creation}
\label{sec:ch11-type-def}

\begin{listing}[H]
\caption{Type definition and blockchain creation}
\label{lst:ch11-create}
\begin{minted}[fontsize=\footnotesize]{rust}
use crate::error::{BtcError, Result};
use crate::primitives::block::Block;
use crate::primitives::blockchain::Blockchain;
use crate::primitives::transaction::Transaction;
use crate::wallet::WalletAddress;
use sled::{Db, Tree};
use std::env;
use std::sync::Arc;
use tokio::sync::RwLock as TokioRwLock;

const DEFAULT_TIP_BLOCK_HASH_KEY: &str = "tip_block_hash";
const DEFAULT_EMPTY_TIP_BLOCK_HASH_VALUE: &str = "empty";
const DEFAULT_BLOCKS_TREE: &str = "blocks1";
const DEFAULT_TREE_DIR: &str = "data1";

#[derive(Clone, Debug)]
pub struct BlockchainFileSystem {
    blockchain: Blockchain<Db>,
    file_system_tree_dir: String,
}

impl BlockchainFileSystem {
    pub async fn create_blockchain(
        genesis_address: &WalletAddress,
    ) -> Result<Self> {
        let blocks_tree_name = env::var("BLOCKS_TREE")
            .unwrap_or_else(|_| DEFAULT_BLOCKS_TREE.to_string());
        let dir = env::var("TREE_DIR")
            .unwrap_or_else(|_| DEFAULT_TREE_DIR.to_string());
        let path = std::env::current_dir()?
            .join(dir.clone());

        let db = sled::open(path)?;
        let blocks_tree = db.open_tree(blocks_tree_name.clone())?;

        // Check if chain exists, else create genesis
        let tip_hash =
            if let Some(data) =
                blocks_tree.get(DEFAULT_TIP_BLOCK_HASH_KEY)? {
            String::from_utf8(data.to_vec())?
        } else {
            // Initialize: create genesis block,
            // persist + set tip atomically
            let coinbase_tx =
                Transaction::new_coinbase_tx(genesis_address)?;
            let block =
                Block::generate_genesis_block(&coinbase_tx);
            Self::update_blocks_tree(&blocks_tree, &block)
                .await?;
            String::from(block.get_hash())
        };

        Ok(BlockchainFileSystem {
            blockchain: Blockchain {
                tip_hash: Arc::new(TokioRwLock::new(tip_hash)),
                db,
                is_empty: false,
            },
            file_system_tree_dir: blocks_tree_name,
        })
    }
}
\end{minted}
\end{listing}

% ──────────────────────────────────────────────────────────────
\section{Open Existing or Empty Blockchain}
\label{sec:ch11-open-blockchain}

\begin{listing}[H]
\caption{Open existing or empty blockchain}
\label{lst:ch11-open}
\begin{minted}[fontsize=\footnotesize]{rust}
impl BlockchainFileSystem {
    pub async fn open_blockchain() -> Result<Self> {
        let blocks_tree_name = env::var("BLOCKS_TREE")
            .unwrap_or_else(|_| DEFAULT_BLOCKS_TREE.to_string());
        let dir = env::var("TREE_DIR")
            .unwrap_or_else(|_| DEFAULT_TREE_DIR.to_string());
        let path = std::env::current_dir()?.join(dir);

        let db = sled::open(path)?;
        let blocks_tree = db.open_tree(blocks_tree_name.clone())?;

        let tip_bytes = blocks_tree
            .get(DEFAULT_TIP_BLOCK_HASH_KEY)?
            .ok_or(BtcError::BlockchainNotFoundError(
                "No blockchain found".into(),
            ))?;
        let tip_hash = String::from_utf8(tip_bytes.to_vec())?;

        Ok(BlockchainFileSystem {
            blockchain: Blockchain {
                tip_hash: Arc::new(TokioRwLock::new(tip_hash)),
                db,
                is_empty: false,
            },
            file_system_tree_dir: blocks_tree_name,
        })
    }

    pub async fn open_blockchain_empty() -> Result<Self> {
        let blocks_tree_name = env::var("BLOCKS_TREE")
            .unwrap_or_else(|_| DEFAULT_BLOCKS_TREE.to_string());
        let dir = env::var("TREE_DIR")
            .unwrap_or_else(|_| DEFAULT_TREE_DIR.to_string());
        let path = std::env::current_dir()?.join(dir);

        let db = sled::open(path)?;
        Ok(BlockchainFileSystem {
            blockchain: Blockchain {
                tip_hash: Arc::new(
                    TokioRwLock::new("empty".to_string())),
                db,
                is_empty: true,
            },
            file_system_tree_dir: blocks_tree_name,
        })
    }
}
\end{minted}
\end{listing}

% ──────────────────────────────────────────────────────────────
\section{Atomic Storage Primitive: \texttt{update\_blocks\_tree}}
\label{sec:ch11-atomic-primitive}

When we create genesis or mine locally, we want ``block insert'' and ``tip pointer update'' to happen as a single atomic unit.

\begin{listing}[htbp]
\caption{Atomic insert + tip update}
\label{lst:ch11-atomic}
\begin{minted}[fontsize=\footnotesize]{rust}
async fn update_blocks_tree(
    blocks_tree: &Tree,
    block: &Block
) -> Result<()> {
    // Hash is the key; serialized block is the value.
    let block_hash = block.get_hash();
    let block_ivec = IVec::try_from(block.clone())?;

    // sled transactions give us atomicity across
    // multiple inserts into the same Tree.
    let transaction_result: TransactionResult<(), ()> =
        blocks_tree.transaction(|tx_db| {
        // 1) Persist the block bytes under its hash.
        let _ = tx_db.insert(
            block_hash,
            block_ivec.clone()
        )?;
        // 2) Move the "tip pointer" to the same hash.
        // If the process crashes after this transaction
        // commits, both updates are present.
        let _ = tx_db.insert(
            DEFAULT_TIP_BLOCK_HASH_KEY,
            block_hash
        )?;
        Ok(())
    });
    transaction_result
        .map(|_| ())
        .map_err(|e|
            BtcError::BlockchainDBconnection(
                format!("{:?}", e)))
}
\end{minted}
\end{listing}

% ──────────────────────────────────────────────────────────────
\section{Tip Pointer API}
\label{sec:ch11-tip-api}

The tip is also cached in memory (an async \code{RwLock<String>}) for fast access during runtime.

\begin{listing}[htbp]
\caption{Tip pointer helpers}
\label{lst:ch11-tip-helpers}
\begin{minted}[fontsize=\footnotesize]{rust}
pub async fn get_tip_hash(&self) -> Result<String> {
    // Read-lock: multiple readers can read
    // the cached tip concurrently.
    let tip_hash = self.blockchain.tip_hash.read().await;
    Ok(tip_hash.clone())
}

async fn set_tip_hash(&self, new_tip_hash: &str)
    -> Result<()> {
    // Write-lock: tip updates are exclusive.
    let mut tip_hash =
        self.blockchain.tip_hash.write().await;
    // Update the in-memory cached tip
    // (DB tip key updates happen elsewhere).
    *tip_hash = String::from(new_tip_hash);
    Ok(())
}
\end{minted}
\end{listing}

% ──────────────────────────────────────────────────────────────
\section{Read Path: Height, Get-by-Hash, and Inventory}
\label{sec:ch11-read-path}

The network layer and UI need:

\begin{itemize}
  \item current height (\code{get\_best\_height})
  \item fetch-by-hash (\code{get\_block})
  \item list hashes (\code{get\_block\_hashes}) for sync inventory
\end{itemize}

\begin{listing}[htbp]
\caption{Read path operations}
\label{lst:ch11-read}
\begin{minted}[fontsize=\footnotesize]{rust}
pub async fn get_best_height(&self) -> Result<usize> {
    if self.is_empty() {
        // An "empty handle" reports height 0
        // (genesis is height 1 in this codebase).
        info!("Blockchain is empty, returning height 0");
        Ok(0)
    } else {
        // Open the blocks tree so we can fetch
        // the tip block bytes.
        let block_tree = self
            .blockchain.db
            .open_tree(self.get_blocks_tree_path())
            .map_err(|e|
                BtcError::OpenBlockchainTreeError(
                    e.to_string()))?;
        // Use the cached tip hash to look up the
        // corresponding block record.
        let tip_block_bytes = block_tree
            .get(self.get_tip_hash().await?)
            .map_err(|e|
                BtcError::GetBlockchainError(
                    e.to_string()))?
            .ok_or(BtcError::GetBlockchainError(
                "tip is invalid".to_string()))?;
        // Deserialize bytes back into a Block
        // so we can read its height.
        let tip_block =
            Block::deserialize(tip_block_bytes.as_ref())?;
        Ok(tip_block.get_height())
    }
}

pub async fn get_block(
    &self,
    block_hash: &[u8]
) -> Result<Option<Block>> {
    // Read-by-hash is a straight key-value lookup.
    let block_tree = self
        .blockchain.db
        .open_tree(self.get_blocks_tree_path())
        .map_err(|e|
            BtcError::OpenBlockchainTreeError(
                e.to_string()))?;
    let block_bytes = block_tree
        .get(block_hash)
        .map_err(|e|
            BtcError::GetBlockchainError(e.to_string()))?;

    if let Some(block_bytes) = block_bytes {
        // Found a record: decode it back into a Block.
        let block =
            Block::deserialize(block_bytes.as_ref())?;
        Ok(Some(block))
    } else {
        // Missing key: caller gets None (not an error).
        Ok(None)
    }
}

pub async fn get_block_hashes(&self)
    -> Result<Vec<Vec<u8>>> {
    // "Inventory" for sync: walk the chain
    // backward and collect hashes.
    let mut iterator = self.iterator().await?;
    let mut blocks = vec![];
    loop {
        match iterator.next() {
            None => break,
            Some(block) => {
                blocks.push(block.get_hash_bytes());
            }
        }
    }
    Ok(blocks)
}
\end{minted}
\end{listing}

% ──────────────────────────────────────────────────────────────
\section{Mining Write Path: \texttt{mine\_block}}
\label{sec:ch11-mining-write}

Local mining creates a new \code{Block}, writes it into sled, updates the in-memory tip, then updates derived state. This is the ``fast path'' for locally-mined blocks --- it bypasses the full consensus mechanism in \code{add\_block} because we know the block is at \code{best\_height + 1} (strictly higher than the current tip).

\begin{listing}[htbp]
\caption{Mining write path}
\label{lst:ch11-mine}
\begin{minted}[fontsize=\footnotesize]{rust}
pub async fn mine_block(
    &self,
    transactions: &[Transaction]
) -> Result<Block> {
    let best_height = self.get_best_height().await?;
    let block = Block::new_block(
        self.get_tip_hash().await?,
        transactions,
        best_height + 1,
    );
    let block_hash = block.get_hash();

    let blocks_tree = self
        .blockchain.db
        .open_tree(self.get_blocks_tree_path())
        .map_err(|e|
            BtcError::BlockchainDBconnection(
                e.to_string()))?;
    Self::update_blocks_tree(&blocks_tree, &block)
        .await?;
    self.set_tip_hash(block_hash).await?;
    self.update_utxo_set(&block).await?;

    Ok(block)
}
\end{minted}
\end{listing}

\begin{notebox}[Stale Mining Protection]
The \code{BlockchainService::mine\_block()} wrapper re-validates transaction inputs under the write lock before calling this method. Between \code{prepare\_mining\_utxo} (which validates with a read lock) and the actual mining (which requires the write lock), a competing block may have been accepted, spending the same inputs. The wrapper catches this race condition and aborts mining with a ``stale mining'' error rather than creating a block with already-spent inputs. This prevents duplicate mining subsidies when multiple miners compete for the same transactions.
\end{notebox}

% ──────────────────────────────────────────────────────────────
\section{Inbound Write Path: \texttt{add\_block}}
\label{sec:ch11-inbound-write}

Inbound blocks are first persisted, then the node runs consensus decisions to determine whether the tip should move or whether a reorganization is required.

\begin{listing}[htbp]
\caption{Inbound block persistence and dedup (part 1)}
\label{lst:ch11-add-block-1}
\begin{minted}[fontsize=\footnotesize]{rust}
pub async fn add_block(&mut self, new_block: &Block)
    -> Result<()> {
    // Inbound: persist first -> then consensus
    // decision -> maybe reorg.
    let block_tree = self.blockchain.db
        .open_tree(self.get_blocks_tree_path())?;

    if self.is_empty() {
        // First block always becomes tip
        self.set_not_empty();
        Self::update_blocks_tree(&block_tree, new_block)
            .await?;
        self.set_tip_hash(new_block.get_hash()).await?;
        self.update_utxo_set(new_block).await?;
        return Ok(());
    }

    // Dedup: skip if already present
    if block_tree.get(new_block.get_hash())?
        .is_some() {
        return Ok(());
    }

    // Atomic insert within transaction
    let block_bytes = new_block.serialize()?;
    let tip_hash = self.get_tip_hash().await?;
    block_tree.transaction(|txn| {
        txn.insert(
            new_block.get_hash(),
            block_bytes.clone()
        )?;
        // Eagerly move tip if block is strictly higher
        if new_block.get_height() >
            Block::deserialize(
                txn.get(tip_hash.clone())?
                    .ok_or(...)?)?
            ).get_height() {
            txn.insert(
                DEFAULT_TIP_BLOCK_HASH_KEY,
                new_block.get_hash()
            )?;
        }
        Ok(())
    })?;
\end{minted}
\end{listing}

After persisting the block, we run consensus rules to decide whether to update the tip or reorganize:

\begin{listing}[htbp]
\caption{Fork choice and reorg decision (part 2)}
\label{lst:ch11-add-block-2}
\begin{minted}[fontsize=\footnotesize]{rust}
    // Fork choice: compare height, work, then tie-break
    let current_tip = self.get_tip_hash().await?;
    let current_height = self.get_best_height().await?;

    match new_block.get_height().cmp(&current_height) {
        Ordering::Greater => {
            // Check if block extends our chain
            // or is on a different branch
            if new_block.get_pre_block_hash()
                == current_tip {
                // Normal case: extends our current chain
                self.set_tip_hash(
                    new_block.get_hash()
                ).await?;
                self.update_utxo_set(new_block).await?;
            } else {
                // FORK: block is on a different branch
                // at a higher height.
                // Must reorganize to properly rollback
                // old branch's UTXO.
                self.reorganize_chain(
                    new_block.get_hash()
                ).await?;
            }
        }
        Ordering::Equal => {
            if new_block.get_pre_block_hash()
                == current_tip {
                return Ok(());
            }
            // Competing tips: compare cumulative work
            let current_work =
                self.get_chain_work(&current_tip).await?;
            let new_work = self.get_chain_work(
                new_block.get_hash()
            ).await?;
            match new_work.cmp(&current_work) {
                Ordering::Greater => {
                    self.reorganize_chain(
                        new_block.get_hash()
                    ).await?
                }
                Ordering::Equal => {
                    if self.accept_new_block_tie_break(
                        new_block,
                        &current_tip
                    ).await? {
                        self.reorganize_chain(
                            new_block.get_hash()
                        ).await?;
                    }
                }
                Ordering::Less => {}
            }
        }
        Ordering::Less => {
            // Block stored in DB but not accepted as tip.
            // Available for future reorganizations when
            // higher blocks on this branch arrive.
        }
    }
    Ok(())
}
\end{minted}
\end{listing}

\begin{notebox}[Important Fork Detection Notes]
\begin{enumerate}
  \item \textbf{\code{Ordering::Greater} fork detection}: The original code assumed higher-height blocks always extended the current chain. With block relay across multi-hop networks, a higher-height block on a \emph{different branch} can arrive. Without checking \code{get\_pre\_block\_hash() == current\_tip}, the old branch's UTXO (including coinbase subsidies) would never be rolled back, creating money from nothing. The fix calls \code{reorganize\_chain()} to properly rollback and apply.

  \item \textbf{\code{Ordering::Less} block retention}: Blocks at lower heights are now kept in the database. They are NOT deleted. This is critical because \code{find\_common\_ancestor()} needs to walk both chain branches during reorganization. If rolled-back blocks were deleted, future reorganizations would fail with ``No common ancestor found.'' This matches Bitcoin Core's behavior where all received blocks are kept in \code{blk*.dat} files.

  \item \textbf{\code{rollback\_to\_block} no longer deletes blocks}: During chain reorganization, only the UTXO set is rolled back and the tip pointer is moved. Blocks remain in the database as non-canonical entries. This ensures \code{find\_common\_ancestor()} and \code{apply\_chain\_from\_ancestor()} can always find the necessary blocks.
\end{enumerate}
\end{notebox}

% ──────────────────────────────────────────────────────────────
\section{Exercises}
\label{sec:ch11-exercises}

\begin{enumerate}
  \item \textbf{Storage Round-Trip Test} --- Write a test that creates a block, stores it using \code{BlockchainFileSystem}, retrieves it, and verifies all fields match the original. Pay attention to serialization: does every field survive the round trip?

  \item \textbf{UTXO Retrieval Performance} --- Examine how UTXO lookups work in the Sled storage layer. If the blockchain has 10,000 transactions with an average of 2 outputs each, how many entries are in the UTXO set? What data structure does Sled use internally that makes lookups efficient?
\end{enumerate}

% ──────────────────────────────────────────────────────────────
\section{Common Errors}
\label{sec:ch11-common-errors}

\begin{notebox}[Troubleshooting: Storage Layer]
\textbf{``sled lock file already exists''} --- Another process (or a crashed previous run) still holds the database lock. Delete the \code{db/\_\_sled\_\_/lock} file manually, or ensure the previous node process has fully exited.

\textbf{Deserialization panic after code changes} --- If you modify the \code{Block} struct and restart without clearing the database, sled will try to deserialize old bytes into the new struct layout. Delete the \code{db/} directory to start fresh, or implement a migration path.

\textbf{``tree not found'' on a fresh database} --- Ensure \code{create\_blockchain} is called before \code{open\_blockchain}. The blocks tree must be explicitly created on first run.
\end{notebox}

% ──────────────────────────────────────────────────────────────
\begin{chaptersummary}
\item We explored how the \code{BlockchainFileSystem} uses the Sled embedded database to persist blocks, chain state, and UTXO data to disk.
\item We implemented the storage operations that allow efficient insertion and retrieval of blockchain data.
\item We examined the storage abstraction that decouples chain logic from persistence details, making it possible to swap storage backends.
\item We understood the critical role of atomic transactions in ensuring consistency across crashes and failures.
\end{chaptersummary}

In the next chapter, we build the networking layer that enables nodes to communicate, propagate blocks, and maintain consensus across a distributed network.
